%% Software and Systems Modeling (SoSyM) manuscript
%% Based on the Springer Nature LaTeX template, version 3.1 (December 2024),
%% kept unmodified in ../template/. sn-jnl.cls and sn-mathphys-num.bst sit next
%% to this file because the submission must include all style files.
%%
%% SoSyM submission guidelines (checked 2026-09-13):
%%   - LaTeX with the [iicol] option; upload the full source and a compiled PDF
%%   - abstract of 150 to 250 words, without undefined abbreviations or citations
%%   - numbered citations in square brackets (sn-mathphys-num)
%%   - decimal headings, at most three levels
%%   - a "Statements and Declarations" section with Competing Interests,
%%     plus a data availability statement
%%   - single-blind review, so author names and repository links can stay
%% The template asks for a single .tex file, so avoid \input.
%%
%% Writing rules for the manuscript: plain words, no overclaiming, short
%% sentences, no pronouns, and no em or en dashes in the prose.
%% Sources: the ecorefs README and evaluation notes, the two BPMN2 repositories,
%% and ../literature-review.md (every reference checked online).

%% Allow line breaks at hyphens inside long URLs in the reference list
\PassOptionsToPackage{hyphens}{url}
\documentclass[iicol,pdflatex,sn-mathphys-num]{sn-jnl}

%%%% Standard packages from the template
\usepackage{graphicx}%
\usepackage{multirow}%
\usepackage{amsmath,amssymb,amsfonts}%
\usepackage{amsthm}%
\usepackage{mathrsfs}%
\usepackage[title]{appendix}%
\usepackage{xcolor}%
\usepackage{textcomp}%
\usepackage{manyfoot}%
\usepackage{booktabs}%
\usepackage{algorithm}%
\usepackage{algorithmicx}%
\usepackage{algpseudocode}%
\usepackage{listings}%

%% Code listings for Java and XMI snippets
\lstset{basicstyle=\ttfamily\footnotesize,columns=fullflexible,
  breaklines=true,frame=single,captionpos=b}

\theoremstyle{thmstylethree}%
\newtheorem{definition}{Definition}%

%% The bibliography style prints web addresses with \burl. The class loads
%% breakurl, but \burl stays undefined under pdfLaTeX, so map \burl to \url.
\providecommand{\burl}[1]{\url{#1}}

%% Drafting aid: remove every \todo before submission
\newcommand{\todo}[1]{\textcolor{red}{[TODO: #1]}}

\raggedbottom

\begin{document}

\title[Decentralized Model Persistence]{Towards Decentralized Model Persistence}

\author*[1]{\fnm{Alfa} \sur{Yohannis}}\email{alfa.ryano@pradita.ac.id}
%% Further authors: \author[1]{\fnm{Given} \sur{Family}}\email{name@example.org}

\affil*[1]{\orgdiv{Department}, \orgname{Pradita University}, \orgaddress{\city{City}, \country{Indonesia}}}

\abstract{Model persistence frameworks such as the Eclipse Modeling Framework (EMF) identify a model by a location, for example a file path. The content at a location can change while the location stays the same. The InterPlanetary File System (IPFS) identifies a file by a content identifier (CID) instead. The CID is computed from the file content, so every change to a model produces a new identifier. The paper examines the consequences for model persistence and for references between model resources. An EMF resource implementation saves models to a local IPFS node and loads the models back by CID. References can point to fixed CIDs or to mutable names of the InterPlanetary Name System (IPNS). A cascade save stores every resource with a reference to a changed resource again. IPFS support was also added to the Eclipse BPMN2 Modeler. A benchmark used generated models of 100 to 10,000 elements. A save through IPFS took 9 to 10~ms longer than a save to a local XML Metadata Interchange (XMI) file. A case study split 39 MoDisco models of the Eclipse Java development tools (1.04~GiB) into 6,933 resources and published the resources on IPFS. For eight components with 1,639 resources, EMF Compare found no differences between the local models and the copies fetched from IPFS.}

\keywords{Model-driven engineering, Model persistence, Eclipse Modeling Framework, IPFS, Content-addressed storage, Cross-resource references}

\maketitle

\section{Introduction}\label{sec:introduction}

Model-driven engineering (MDE) uses models to manage the complexity of platforms and to express domain concepts~\cite{schmidt2006mde}. MDE tools store models in files or in model repositories. The Eclipse Modeling Framework (EMF) stores each model in a resource, and a Uniform Resource Identifier (URI) names the resource~\cite{steinberg2008emf}. Resources are usually serialized in the XML Metadata Interchange (XMI) format~\cite{omg2015xmi}. Common URIs such as \texttt{file://} and \texttt{platform://} name a location. The content at a location can change, but the URI stays the same. A reference to another resource therefore resolves to the current content at the location. The referenced model may have changed since the reference was created.

Content-addressed storage takes a different approach. The InterPlanetary File System (IPFS) identifies content by a content identifier (CID)~\cite{benet2014ipfs,ipfsdocscid}. The CID is computed from the content, so a changed file receives a new CID~\cite{ipfsdocsipns}. Content under a given CID cannot change. Any IPFS node with a copy of the content can serve the content~\cite{trautwein2022ipfs}. The InterPlanetary Name System (IPNS) adds mutable names on top of CIDs~\cite{ipfsdocsipns}.

Content addressing has two consequences for models. First, a reference to a CID pins the exact version of the referenced model. Package managers use a similar idea. A lock file records the exact dependency tree of a project~\cite{npmpackagelock}, and Nix derives unique component paths from cryptographic hashes~\cite{dolstra2004nix}. Second, every save of a model creates a new CID. References to the model then need a deliberate update, and tools need support for the update. Dependency updates are a known source of problems in software ecosystems~\cite{cox2019dependencies,decan2019ecosystems}.

Research on model persistence has mainly addressed scalability~\cite{kolovos2013roadmap}. Database-backed solutions load parts of a large model on demand~\cite{pagan2015repository,daniel2017neoemf}. Model versioning systems track changes between versions of a model~\cite{koegel2010emfstore,altmanninger2009survey}. None of the reviewed approaches stores EMF models on content-addressed storage (Section~\ref{sec:related}).

The paper studies content-addressed persistence for EMF with three research questions:

%% Research questions adapted from the ecorefs README.md ("Research questions").
%% The README's "Publication plan" targets RQ2 at MODELS and RQ1/RQ3 at the
%% journal version; drop or shorten RQ2 here if RQ2 is published separately.
\begin{enumerate}
\item[RQ1] What are the architectural challenges of extending model persistence frameworks to support content-addressed decentralized storage, and how can the challenges be addressed?
\item[RQ2] How can cross-resource references be maintained and resolved in an immutable, content-addressed environment where resource identifiers change on every modification?
\item[RQ3] What is the performance overhead of content-addressed model persistence, and under what conditions is content-addressed persistence practical for MDE toolchains?
\end{enumerate}

The paper makes four contributions:
\begin{itemize}
\item An EMF resource implementation for IPFS. The implementation saves models as XMI files on a local IPFS node and loads models by CID or IPNS name (Section~\ref{sec:rq1}).
\item A description of four strategies for cross-resource references on IPFS and a cascade save for references by CID (Section~\ref{sec:rq2}).
\item IPFS support in the Eclipse BPMN2 Modeler and in the generated BPMN2 tree editor (Section~\ref{sec:tool}).
\item An evaluation with generated models and with large MoDisco models of Java code (Section~\ref{sec:evaluation}).
\end{itemize}

Section~\ref{sec:background} introduces EMF persistence and IPFS. Section~\ref{sec:rq1} describes the persistence layer, and Section~\ref{sec:rq2} covers references between resources. Section~\ref{sec:tool} presents the tool support. Section~\ref{sec:evaluation} reports the evaluation. Section~\ref{sec:discussion} answers the research questions and discusses threats to validity. Section~\ref{sec:related} covers related work, and Section~\ref{sec:conclusion} concludes the paper.

\section{Background}\label{sec:background}

\subsection{Model persistence in EMF}\label{sec:background-emf}

EMF defines metamodels with Ecore. Ecore is aligned with the Essential MOF subset of the Meta Object Facility~\cite{steinberg2008emf,omg2016mof}. A \texttt{Resource} holds the root objects of a model, and a URI identifies the resource. A \texttt{ResourceSet} groups resources and resolves references between the resources. A resource factory registry maps a URI scheme or a file extension to a resource factory. A \texttt{URIConverter} opens input and output streams for a URI through a list of URI handlers~\cite{steinberg2008emf}.

A reference to an object in another resource is serialized as a URI with a fragment, for example \verb|href="model.xmi#//Person"|. During loading, EMF creates a proxy object for such a reference. The proxy is resolved on first access, and the target resource is loaded when needed.

\subsection{Content-addressed storage with IPFS}\label{sec:background-ipfs}

Content addressing identifies data by a hash of the data. Merkle introduced hash trees for authentication~\cite{merkle1988signature}. Venti used a hash of the block contents as the identifier of a block in archival storage~\cite{quinlan2002venti}. Git stores every object under the SHA-1 checksum of the object content~\cite{chacon2014progit}.

IPFS combines content addressing with a peer-to-peer network~\cite{benet2014ipfs}. Content is stored in blocks, and links between blocks form a Merkle directed acyclic graph (DAG)~\cite{benet2014ipfs}. Nodes locate content through a distributed hash table (DHT) based on Kademlia~\cite{maymounkov2002kademlia,psaras2020ipfs}. Studies describe the design and performance of IPFS~\cite{trautwein2022ipfs} and the size of the public network~\cite{henningsen2020mapping}. A survey compares IPFS with other peer-to-peer data networks~\cite{daniel2022ipfsfriends}.

A directory on IPFS also receives a CID. The directory CID covers all files in the directory. A relative path inside the directory therefore always leads to the same file under the same directory CID. Content stays available only while at least one IPFS node stores the content.

IPNS provides mutable names~\cite{ipfsdocsipns}. An IPNS name is the hash of a public key. The holder of the private key can publish a new CID under the name at any time. Resolving an IPNS name needs an extra lookup before the content can be fetched.

Kubo is an IPFS implementation in Go~\cite{kubo}. A Kubo node offers a Hypertext Transfer Protocol (HTTP) application programming interface (API). The API includes commands such as \texttt{add}, \texttt{cat}, \texttt{name/publish} and \texttt{name/resolve}.

\section{Content-addressed model persistence}\label{sec:rq1}

The ecorefs core plugin adds IPFS persistence to EMF. The plugin is a Java library built with Maven for Java~11. The plugin uses the EMF Ecore library (version 2.23.0) and a Java client for the IPFS HTTP API (version 1.4.3)~\cite{javaipfsclient}.

\subsection{Architecture}\label{sec:rq1-architecture}

Table~\ref{tab:classes} lists the main classes of the plugin. Listing~\ref{lst:setup} shows the configuration of a resource set. The factory is registered for the \texttt{ipfs} and \texttt{ipns} schemes. The URI handler is added at the front of the handler list, so the handler receives IPFS URIs before the default handlers.

\begin{table*}[t]
\caption{Main classes of the ecorefs core plugin}\label{tab:classes}
\footnotesize
\begin{tabular}{@{}p{0.26\textwidth}p{0.68\textwidth}@{}}
\toprule
Class & Role \\
\midrule
\texttt{IPFSResourceImpl} & Subclass of \texttt{XMIResourceImpl}. Saves the resource to IPFS and replaces the resource URI with the new CID or IPNS name. \\
\texttt{IPFSResourceFactoryImpl} & Creates \texttt{IPFSResourceImpl} instances. \\
\texttt{IPFSURIHandlerImpl} & Reads \texttt{ipfs://}, \texttt{ipns://}, \texttt{http://} and \texttt{https://} URIs. \\
\texttt{IPFSModelPersister} & Provides the cascade save (Section~\ref{sec:rq2-cascade}). \\
\texttt{ReferenceMode} & Selects references by CID or by IPNS name through a save option. \\
\botrule
\end{tabular}
\end{table*}

\begin{lstlisting}[language=Java,caption={Configuration of a resource set for IPFS, adapted from the test code},label={lst:setup}]
IPFS ipfs = new IPFS("/ip4/127.0.0.1/tcp/5001");
ResourceSet resourceSet = new ResourceSetImpl();
IPFSResourceFactoryImpl factory =
    new IPFSResourceFactoryImpl(ipfs);
Map<String, Object> protocols = resourceSet
    .getResourceFactoryRegistry()
    .getProtocolToFactoryMap();
protocols.put("ipfs", factory);
protocols.put("ipns", factory);
resourceSet.getURIConverter().getURIHandlers()
    .add(0, new IPFSURIHandlerImpl(ipfs));
\end{lstlisting}

\subsection{Saving a model}\label{sec:rq1-save}

Saving an \texttt{IPFSResourceImpl} takes three steps. First, the resource writes the XMI serialization to a temporary file. The temporary file avoids keeping a large serialized model in memory. Second, the file is added to the local Kubo node, and Kubo returns the CID. Third, the resource URI is replaced with \texttt{ipfs://} followed by the CID.

In IPNS mode, the third step differs. The resource derives a key name from the resource URI. An Ed25519 key is created when no key with the name exists. The new CID is then published under the key, and the resource URI becomes \texttt{ipns://} followed by the key identifier. The key name is kept after the first save, so later saves publish under the same name.

\subsection{Loading a model}\label{sec:rq1-load}

\texttt{IPFSURIHandlerImpl} reads models for the resource set. For an \texttt{ipfs://} URI, the handler removes the fragment and fetches the content with the Kubo \texttt{cat} command. For an \texttt{ipns://} URI, the handler first resolves the name to a CID. The handler also reads plain \texttt{http://} and \texttt{https://} URIs. Output streams for all four schemes are rejected. A model is written only through \texttt{IPFSResourceImpl.save}, because content under an existing CID cannot change.

\subsection{Differences from location-based persistence}\label{sec:rq1-differences}

The implementation shows four points where content addressing departs from the assumptions of the EMF resource model:
\begin{enumerate}
\item A new resource needs a URI before the first save, but the CID exists only after the save. The tests and examples use placeholder URIs such as \texttt{ipfs://pending}.
\item Every save in CID mode assigns a new URI to the resource. Code with a stored copy of the old URI still refers to the old version.
\item An output stream for an existing CID has no meaning. The save operation therefore writes a temporary file and adds the file to IPFS without the output streams of the URI converter.
\item File resources and IPFS resources can use one resource set configuration, because the IPFS handler accepts only IPFS and HTTP URIs. The BPMN example program loads local \texttt{.bpmn} files and IPFS resources with the same configuration.
\end{enumerate}

\section{Cross-resource references}\label{sec:rq2}

\subsection{Version pinning by content addressing}\label{sec:rq2-pinning}

A reference to a resource on IPFS contains the CID of the target resource, for example \verb|href="ipfs://QmB...#//Person"|. The CID fixes the content of the target. When the target model changes, the changed model receives a new CID. The existing reference still resolves to the old version, as long as an IPFS node keeps the old content.

For models developed independently, the behavior pins a dependency in the same way as a lock file~\cite{npmpackagelock,gomodules}. A model keeps working with the version used during validation. The adoption of a newer version becomes an explicit step.

\subsection{Reference management strategies}\label{sec:rq2-strategies}

Table~\ref{tab:strategies} summarizes four strategies for references between resources on IPFS. The core plugin supports direct CID references and IPNS references through the reference mode. The publication tools for partitioned models combine directories, manifests and one IPNS name (Section~\ref{sec:rq2-partitioned}).

\begin{table*}[t]
\caption{Strategies for cross-resource references on IPFS}\label{tab:strategies}
\footnotesize
\begin{tabular}{@{}p{0.16\textwidth}p{0.3\textwidth}p{0.22\textwidth}p{0.24\textwidth}@{}}
\toprule
Strategy & Mechanism & Benefit & Drawback \\
\midrule
Direct CID references & The reference contains the CID of the target & Full reproducibility & Upgrades are manual \\
IPFS directories & Relative paths inside one directory with one root CID & Suits resources developed together & Any change creates a new root CID \\
IPNS indirection & The reference contains a mutable IPNS name & Resolves to the latest version & No immutability, and resolution needs an extra lookup \\
Manifest or registry & A document maps logical names to CIDs & Flexible & The manifest also needs management \\
\botrule
\end{tabular}
\end{table*}

\subsection{Cascade save}\label{sec:rq2-cascade}

In CID mode, a change to one resource requires new saves of every resource with a reference to the changed resource. \texttt{IPFSModelPersister.cascadeSave} automates the saves (Algorithm~\ref{alg:cascade}).

\begin{algorithm}
\caption{Cascade save in CID mode ($S$ holds the resources saved so far)}\label{alg:cascade}
\begin{algorithmic}[1]
\Procedure{CascadeSave}{$r, S$}
\If{$r \in S$}
\State \Return
\EndIf
\State save $r$ to IPFS
\State $S \gets S \cup \{r\}$
\State $X \gets$ all cross-references in the resource set
\State $P \gets$ resources with a reference in $X$ into $r$
\ForAll{$p \in P \setminus S$}
\State \Call{CascadeSave}{$p, S$}
\EndFor
\EndProcedure
\end{algorithmic}
\end{algorithm}

The procedure first saves the changed resource, and the save updates the URI of the resource to the new CID. The EMF cross referencer then collects all cross-references in the resource set. Each resource with a reference into the saved resource is saved in the same way. During serialization, EMF writes the current URI of the target into each reference. The rewritten references therefore point to the new CID without any text replacement. The procedure resembles a Git commit, where a changed file leads to new tree objects up to the root~\cite{chacon2014progit}. The set of saved resources stops the recursion for cyclic references.

In IPNS mode, \texttt{cascadeSave} saves only the changed resource. References by IPNS name resolve to the newly published CID, so the referencing resources need no new save.

Cycles need care. Resource A may refer to resource B while B refers to A. In such a cycle, one direction always refers to an outdated CID. A new save of A would change the CID of A again, so the cycle cannot settle. \texttt{IPFSModelPublishTool} therefore checks a partitioned model for cycles before a publication in CID mode. The tool stops with an error when the resource graph contains a cycle.

\todo{Check resources with two reference paths to the changed resource. In the case study, collaboration.bpmn refers to both processes, and the saved-set check may skip a needed second save of the collaboration. A save in reverse topological order would avoid the problem.}

\subsection{Publishing partitioned models}\label{sec:rq2-partitioned}

Large models are often split into many resources. The MoDisco case study (Section~\ref{sec:evaluation-modisco}) uses one root resource and one resource per Java compilation unit. \texttt{IPFSVersionManifestPublisher} publishes such a model in four layers:
\begin{enumerate}
\item Each compilation-unit resource is an immutable file inside the directory of a component.
\item Each plug-in or module forms a component. The component directory is published as one IPFS directory with a component root CID. A component manifest lists the resources of the component and the dependencies of the component.
\item A project version manifest records the CID of the root model and the CIDs of shared resources. For each component, the manifest records the component root CID and the CID of the component manifest.
\item One IPNS name points to the current project version manifest. A local file records the IPNS name and the CID of the current manifest.
\end{enumerate}

The layers combine three strategies from Table~\ref{tab:strategies}: directories, a manifest and IPNS indirection. A new publication after a change to one component gives the changed component a new root CID. The other components keep the same root CIDs, because the content of the other components is unchanged. The new project version manifest is then published under the same IPNS name.

\section{Tool support}\label{sec:tool}

\subsection{Eclipse BPMN2 Modeler}\label{sec:tool-modeler}

The Eclipse BPMN2 Modeler is a graphical tool for creating and editing Business Process Model and Notation (BPMN) diagrams~\cite{bpmn2modeler,omg2010bpmn}. The modeler builds the diagram editors with Graphiti~\cite{graphiti}. IPFS support was added to version 1.5.4 of the modeler, and the modified plug-ins target Java~17. A menu named BPMN2 IPFS offers four commands (Table~\ref{tab:commands}).

\begin{table*}[t]
\caption{IPFS commands in the Eclipse BPMN2 Modeler}\label{tab:commands}
\footnotesize
\begin{tabular}{@{}p{0.26\textwidth}p{0.68\textwidth}@{}}
\toprule
Command & Behavior \\
\midrule
Open BPMN2 Model from IPFS & Accepts a CID, an \texttt{/ipfs/} path, an \texttt{ipfs://} URI or an IPNS reference. Fetches the BPMN file through Kubo, stores the file in the workspace and opens the file in the BPMN2 editor. \\
Publish Current BPMN2 Model to IPFS & Saves the editor if needed and uploads the \texttt{.bpmn} file. Shows the CID, or publishes the CID under the configured IPNS name. \\
Publish BPMN2 Project to IPFS & Publishes every \texttt{.bpmn} and \texttt{.bpmn2} file of the project by CID. Publishes a manifest with the root model and the CID of each relative path, optionally under an IPNS name. \\
Open BPMN2 Project from IPFS & Downloads a project manifest, restores the listed files into a new workspace folder and opens the root model. \\
\botrule
\end{tabular}
\end{table*}

The project manifest is a JavaScript Object Notation (JSON) file. Each project stores four settings on the BPMN2 property page. The settings are the Kubo API URL, a default load reference, the publish mode (CID or IPNS) and a default IPNS key name. The modeler core also registers an IPFS URI handler in the modeler resource set. References with \texttt{ipfs://} inside BPMN files then load through Kubo.

The integration publishes the BPMN model file, and the Graphiti \texttt{.diagram} file stays local. IPNS publishing needs a key name that already exists in the local Kubo node. Without an active workspace folder, a model opened from IPFS is stored in a temporary file.

\subsection{Generated BPMN2 tree editor}\label{sec:tool-editor}

The BPMN~2.0 metamodel of Eclipse includes a tree editor generated by EMF. The generated editor received a similar URI handler. The handler reads \texttt{ipfs://} and \texttt{ipns://} URIs through Kubo and rejects writes. The Kubo API URL comes from a system property, an environment variable or the default address \texttt{http://127.0.0.1:5001}.

The metamodel repository with the handler builds with Maven and Tycho against Eclipse 2022-06 on Java~17. The structural self-tests of the modeler extension pass. \todo{Build the modeler with Tycho. The BPMN2 update site in the build file of the modeler is no longer available.}

\section{Evaluation}\label{sec:evaluation}

The evaluation uses the tests and example programs of the ecorefs repository for RQ1 to RQ3. A case study with large models complements the tests (Section~\ref{sec:evaluation-modisco}).

\subsection{Setup}\label{sec:evaluation-setup}

The benchmarks use a local Kubo node in Docker on port 5001. The tests of the core plugin run with JUnit~5. Each benchmark first runs three warm-up iterations and then ten measured iterations. The tests measure wall-clock time with \texttt{System.nanoTime}. The test code truncates each measurement and each mean to whole milliseconds. \todo{Hardware, operating system, JDK version, Kubo version and Docker settings.}

\todo{Springer asks for a statement on any use of large language models beyond copy editing, placed in the methods part.}

\subsection{RQ1: Functional validation}\label{sec:evaluation-rq1}

An example program (\texttt{BpmnIpfsExample}) exercises the persistence layer with two BPMN models. The program saves both models to IPFS and loads the first model again by CID. The program then resolves an element of the first model through a URI fragment. The second model refers to the process of the first model with an \texttt{ipfs://} URI and a fragment.

The self-tests of the BPMN2 Modeler extension check reference parsing, command registration, bundle metadata and the documentation. An optional mode of the self-tests publishes a BPMN file to Kubo and downloads the file again. \todo{Report the output of the example program and of the self-tests with a running Kubo node.}

\subsection{RQ2: Reference evolution}\label{sec:evaluation-rq2}

The cascade test builds two resources with one object each. The object in resource A has a non-containment reference to the object in resource B. The test calls \texttt{cascadeSave} on B and measures the time for both saves. Over ten runs, the cascade save took 15~ms on average, with a median of 15~ms and a standard deviation of 2.28~ms. The measurement covers two IPFS saves of very small resources. The test does not measure the cross-reference search separately, and the test does not check the rewritten reference. \todo{Add an assertion on the rewritten reference and measure longer chains.}

A BPMN case study prepares a longer chain of three resources. \texttt{payment-process.bpmn} defines a payment process with three tasks. \texttt{order-process.bpmn} refers to the payment process. \texttt{collaboration.bpmn} refers to the order process and to the payment process. A change to the payment process should therefore give new CIDs to all three files. \todo{Run the case study with real CIDs and report the CIDs before and after the cascade save.}

\subsection{RQ3: Performance}\label{sec:evaluation-rq3}

The benchmark uses generated models with 100, 1,000 and 10,000 elements. The models conform to a small metamodel with one class, \texttt{Component}. Each component has an integer identifier and contains child components. In each iteration, the test copies the model into two new resources. One copy is saved as an XMI file, and the other copy is saved to IPFS. The IPFS copy is then loaded again through the new CID. Table~\ref{tab:rq3} shows the results.

\begin{table*}[t]
\caption{Save and load times for generated models over ten measured runs, given as mean / median / standard deviation in ms. XMI load times are not reported (Section~\ref{sec:evaluation-rq3}).}\label{tab:rq3}
\begin{tabular*}{\textwidth}{@{\extracolsep\fill}lccc@{}}
\toprule
Elements & XMI save & IPFS save & IPFS load \\
\midrule
100 & 1 / 1 / 0.67 & 10 / 10 / 0.70 & 3 / 3 / 1.02 \\
1,000 & 4 / 4 / 0.66 & 13 / 13 / 1.11 & 7 / 7 / 1.25 \\
10,000 & 31 / 33 / 4.34 & 41 / 38 / 11.15 & 39 / 40 / 7.30 \\
\botrule
\end{tabular*}
\end{table*}

Saving to IPFS took 9 to 10~ms longer than saving to an XMI file at all three sizes. The difference changed little with model size. The small change suggests a fixed cost per call to the local Kubo node. The standard deviation of IPFS saves grew to 11.15~ms at 10,000 elements. Loading from IPFS took 3~ms, 7~ms and 39~ms on average. A separate shorter run gave similar values for 100 and 1,000 elements. IPFS saves took 12~ms and 13~ms, and IPFS loads took 4~ms and 7~ms.

Table~\ref{tab:rq3} leaves out the XMI load times. In the current test, the saved XMI resource stays in the resource set. The load step therefore returns the resource from memory without reading the file. \todo{Load the XMI file in a fresh resource set and repeat the benchmark.}

The benchmark covers models up to 10,000 elements and a local Kubo node. Network latency between IPFS nodes is outside the measurement. \todo{Add larger models, 30 runs per size, a network file system baseline and a statistical test, as planned in the project README. The benchmark code already lists 100,000 and 500,000 elements.}

\subsection{Case study: MoDisco models of Eclipse JDT}\label{sec:evaluation-modisco}

The case study uses Java models of four repositories of the Eclipse Java development tools (JDT). MoDisco reverse engineers software systems into models with MDE techniques~\cite{bruneliere2010modisco,bruneliere2014modisco}. A headless Eclipse workspace imported the Java projects of each repository. The MoDisco Java discoverer then produced one model per Java project. Table~\ref{tab:dataset} summarizes the resulting 39 XMI files with a total size of 1.04~GiB. One project of \texttt{eclipse.jdt.ui} was skipped and recorded in a failure report.

\begin{table*}[t]
\caption{MoDisco models of Eclipse JDT. Elements are XML elements in the XMI files, and cross-reference targets are serialized references to other elements.}\label{tab:dataset}
\begin{tabular*}{\textwidth}{@{\extracolsep\fill}lrrrr@{}}
\toprule
Repository & Files & Size (MiB) & Elements & Cross-reference targets \\
\midrule
\texttt{eclipse.jdt.core} & 1 & 289.67 & 1,159,412 & 1,661,288 \\
\texttt{eclipse.jdt.core.binaries} & 1 & 0.00 & 10 & 0 \\
\texttt{eclipse.jdt.debug} & 9 & 263.56 & 1,001,753 & 1,522,968 \\
\texttt{eclipse.jdt.ui} & 28 & 516.44 & 2,121,557 & 3,130,907 \\
\midrule
Total & 39 & 1,069.67 & 4,282,732 & 6,315,163 \\
\botrule
\end{tabular*}
\end{table*}

The 39 models were merged into one model. A partitioner then split the merged model into one root resource and one resource per compilation unit. The partitioned model has 6,933 resources. The resources belong to 26 components, one component per plug-in or module. The components share 4 resources, and the component graph has 92 dependency edges.

The merged model was published on IPFS as one file. The partitioned model was published as one IPFS directory and through the layers of Section~\ref{sec:rq2-partitioned}.

The case study uses four kinds of checks. Byte-level checks compare the published merged model, the project version manifest and the partitioned root model with the local files. Loader tests load the merged model and the partitioned root model, navigate into the models and resolve cross-references. A traversal test visits every element of the logical model. In the traversal test, the merged model and the partitioned resource set both contained 4,282,694 elements. \todo{Attach the test reports and the output of the byte-level checks.}

The fourth check compares selected components with copies fetched from IPFS. A script downloads the published root model and the directories of the selected components. The root model and the component manifests are compared byte by byte. EMF Compare compares the XMI resources of the selected components~\cite{emfcompare}. Shared auxiliary resources and unselected components are excluded, because the shared resources would dominate every comparison of a single component. The byte-level checks cover the shared resources instead. Table~\ref{tab:compare} lists four comparison runs.

\begin{table}[t]
\caption{EMF Compare runs on selections of components. Size is the XMI size of one side.}\label{tab:compare}
\begin{tabular}{@{}rrrr@{}}
\toprule
Components & Resources & Size (MiB) & Time (ms) \\
\midrule
4 & 44 & 3.63 & 6,181 \\
6 & 223 & 20.75 & 7,138 \\
8 & 1,639 & 240.61 & 168,992 \\
8 (recheck) & 1,639 & 240.61 & 186,696 \\
\botrule
\end{tabular}
\end{table}

In the recheck on 24 March 2026, the comparison of eight components with 1,639 resources finished without differences in 186,696~ms. The component \texttt{org.eclipse.jdt.core} dominated the comparison time. The component has 1,066 resources with 191.62~MiB of XMI. In the earlier run of the same selection, the component took 152,315~ms of 168,992~ms (about 90\%). The time for the same component varied between runs, for example from 519~ms to 1,385~ms for the smallest component.

According to the notes of the case study, whole-model scans ran faster on the merged model. The partitioned model loaded faster at the start and suited checks of single components. \todo{Add the measured load and traversal times.}

\section{Discussion}\label{sec:discussion}

\subsection{Answers to the research questions}\label{sec:answers}

For RQ1, the main challenges come from the changing identity of a resource. A resource has no final URI before the first save, the URI changes with every save, and content under an existing CID cannot be overwritten. The implementation handles the challenges with placeholder URIs and with a new URI after each save. A read-only URI handler covers loading. In the example program, loading by CID and navigation by URI fragment work through the standard resource set API. \todo{Test proxy resolution across ipfs:// references.}

For RQ2, references by CID pin versions. A deliberate upgrade requires new saves along the chain of referencing resources. The cascade save automates the saves for acyclic references. IPNS names avoid the cascade but give up immutability. For the large partitioned model, component directories, manifests and one IPNS name gave a workable publication structure.

For RQ3, a save through a local Kubo node added 9 to 10~ms per model for models up to 10,000 elements. Loading from IPFS took up to 39~ms. The measured costs appear acceptable for interactive editing of models up to 10,000 elements. Larger models, remote nodes and IPNS resolution still need measurement.

\subsection{Limitations}\label{sec:limitations}

The approach loads each resource as a whole. Database-backed persistence can load parts of a model on demand~\cite{pagan2011morsa,daniel2017neoemf}, but the IPFS layer cannot. The approach also offers no query support comparable to OCL queries inside a database~\cite{omg2014ocl,daniel2016mogwai}. IPNS resolution needs an extra lookup and can take longer than direct CID access. In CID mode, a change to a leaf resource causes saves along the whole dependency chain. Cyclic references always leave one outdated CID. Published models stay available only while at least one node stores the content.

\subsection{Threats to validity}\label{sec:threats}

The main threat to internal validity is the XMI load measurement, because the test did not read the XMI file (Section~\ref{sec:evaluation-rq3}). The truncation to whole milliseconds also hides differences below 1~ms. The cascade test covers only two very small resources.

External validity is limited by the setup. The benchmarks use one generated metamodel and a local Kubo node. The case study covers Java code models only. Results may differ for other metamodels, for remote IPFS nodes and for the public IPFS network.

For construct validity, the wall-clock time of single save and load calls stands for the cost of persistence. The EMF Compare check covered eight components and left out the shared resources.

\section{Related work}\label{sec:related}

\subsection{Scalable model persistence}\label{sec:related-persistence}

Kolovos et al. propose a research roadmap for scalability in MDE, covering large models, collaborative modeling, and scalable queries and transformations~\cite{kolovos2013roadmap}. Morsa stores large models in a NoSQL database and loads model parts on demand~\cite{pagan2011morsa}. The journal version adds incremental store to the Morsa repository~\cite{pagan2015repository}. Neo4EMF provides a scalable persistence layer for EMF~\cite{benelallam2014neo4emf}. NeoEMF offers several database back ends in one framework, because a single data store rarely suits all modeling activities~\cite{daniel2017neoemf}. Mogwaï translates OCL queries to Gremlin and runs the queries inside the database~\cite{daniel2016mogwai}. Scheidgen compares model representations such as XMI trees and relational data in CDO~\cite{scheidgen2013reference}. The CDO model repository stores and versions models in a central repository with relational, NoSQL or in-memory databases~\cite{cdo}. The approaches identify model elements inside a store. The identifiers stay stable while the content changes.

\subsection{Model repositories and versioning}\label{sec:related-versioning}

ReMoDD collects model-driven development artifacts for research and education~\cite{france2007remodd}. Di Rocco et al. discuss model repositories for collaborative modeling, tool interoperability and model reuse~\cite{dirocco2015repositories}. EMFStore versions EMF models with operation-based change tracking, conflict detection and merging~\cite{koegel2010emfstore}. Altmanninger et al. survey version control systems for models with a focus on three-way merging~\cite{altmanninger2009survey}. Brosch et al. introduce the research challenges of model versioning~\cite{brosch2012versioning}. Model comparison supports model composition and transformation testing~\cite{kolovos2006comparison}, and EMF Compare compares and merges models of any metamodel~\cite{emfcompare}. Hawk indexes large collections of models kept in version control systems~\cite{barmpis2013hawk}. A mapping study classifies approaches to collaborative model-driven software engineering~\cite{franzago2018collaborative}. The versioning approaches record or compute differences between versions. Content addressing instead gives each version a fixed identifier derived from the content.

\subsection{Content-addressed storage outside MDE}\label{sec:related-content}

Content addressing is established in several areas. Venti enforced write-once archival storage with hash-based block identifiers~\cite{quinlan2002venti}. Git is a content-addressable file system~\cite{chacon2014progit}. Nix computes unique component paths from cryptographic hashes, so several versions of a component can coexist~\cite{dolstra2004nix}. Software Heritage treats software source code as a digital object with a need for preservation~\cite{dicosmo2017heritage}. On IPFS, InterPlanetary Wayback stores and replays web archives~\cite{kelly2016wayback}. Other work combines IPFS with blockchains for document version control~\cite{nizamuddin2019document} and for a peer-to-peer file system with content service providers~\cite{chen2017improved}. A survey compares IPFS with other decentralized data networks~\cite{daniel2022ipfsfriends}. None of the reviewed approaches applies content addressing to EMF models or to references between model resources.

\subsection{Dependency management and reproducibility}\label{sec:related-dependencies}

Package managers record exact dependency versions. The npm lock file describes the exact dependency tree, so later installations can create identical trees~\cite{npmpackagelock}. Go modules record dependency versions and checksums~\cite{gomodules}. Studies show the risks of dependency updates in software ecosystems~\cite{cox2019dependencies,decan2019ecosystems}. Reproducible builds check whether binaries correspond to the source code~\cite{lamb2022reproducible}. Version pinning by CID applies the same idea to models, and a cascade save plays the role of a deliberate dependency upgrade.

\section{Conclusion}\label{sec:conclusion}

The paper studied EMF model persistence on IPFS. An EMF resource implementation saves models to IPFS and assigns the new CID to the resource URI. A URI handler loads models by CID or IPNS name. References by CID pin exact versions, and a cascade save updates referencing resources in CID mode. IPFS support was also added to the Eclipse BPMN2 Modeler and to the generated BPMN2 tree editor.

With a local Kubo node, a save through IPFS added 9 to 10~ms per model for models up to 10,000 elements. The MoDisco case study published 1.04~GiB of Java models as 6,933 resources. EMF Compare found no differences for eight components with 1,639 resources.

Future work includes a corrected XMI load benchmark, larger models and remote IPFS nodes. The cascade save also needs tests with longer reference chains and with several reference paths to one resource.

\backmatter

\bmhead{Acknowledgements}

\todo{Optional.}

\section*{Statements and Declarations}

\bmhead{Funding}
\todo{Funding statement.}

\bmhead{Competing interests}
\todo{Competing interests statement.}

\bmhead{Ethics approval and consent to participate}
\todo{Write ``Not applicable'' if the item does not apply.}

\bmhead{Consent for publication}
\todo{Write ``Not applicable'' if the item does not apply.}

\bmhead{Data availability}
The benchmark summaries are in the \texttt{eval/results} folder of the ecorefs repository. \todo{Archive the MoDisco dataset (1.04~GiB) and the published CIDs with a DOI.}

\bmhead{Code availability}
The implementation is available at
\url{https://github.com/alfa-yohannis/ecorefs}, the BPMN2 Modeler integration at
\url{https://github.com/alfa-yohannis/org.eclipse.bpmn2-modeler-1.5.4-ecorefs},
and the BPMN~2.0 metamodel changes at
\url{https://github.com/alfa-yohannis/org.eclipse.bpmn2-ecorefs}.

\bmhead{Author contribution}
\todo{Contribution of each author.}

\bibliography{references}

\end{document}
